\documentclass[10pt]{article}
\usepackage[T1]{fontenc}
\usepackage{lmodern}
\usepackage[margin=0.8in]{geometry}
\usepackage{enumitem}
\setlist{itemsep=1pt,topsep=2pt,parsep=0pt}
\usepackage{amsmath,amssymb,amsthm}
\usepackage{booktabs}
\usepackage[hidelinks,pdftitle={Day 224 class-4 escalation, correction, replies (Rick to Clio)},pdfauthor={Rick}]{hyperref}
\newcommand{\WIPHASH}{\texttt{61538d2}}
\newcommand{\file}[1]{{\footnotesize\texttt{\detokenize{#1}}}}
\newcommand{\lin}{\mathrm{lin}_e}
\newcommand{\grade}[1]{\textsf{[#1]}}
\newcommand{\Lead}{\mathrm{Lead}}
\newtheorem{theorem}{Theorem}[section]

\title{Day 224: Box Complement, block multiplicativity, and the class-4 wall (escalation);\\
a correction to my previous note; the block law at general $q$}
\author{Rick}
\date{2026-10-06}

\begin{document}
\sloppy\raggedright
\setlength{\parindent}{0pt}\setlength{\parskip}{3pt}
\maketitle
\vspace{-1.5em}
\noindent\textbf{Author:} Rick. \textbf{Date:} 2026-10-06. \textbf{Recipient:} Clio Vega.\\
\textbf{WIP commit:} \WIPHASH{} of \texttt{grandpa-rick/work-in-progress} (the commit that contains this source file).
The math source is \file{proofs/2026-10-06-day224-Gprime-second-order.md}, pushed earlier at \texttt{2de48f3}.
The escalation data was added at \texttt{dca8e76}.

\paragraph{Grades.} Every bracketed grade is copied from my registry
(\file{registry/hikita-star-dominance-support.json}). The scale is
hunch $<$ sketched $<$ computed $<$ checked-sober $<$ proved. All ``proved'' grades are my own, and none of the
Day~224 nodes has been peer-reviewed. Where a claim has no registry node, I say so and give the grade from the proof file.

\paragraph{Contents.} \S1: Day 224 results, then an escalation of the one open case (``class~4''), with an attack plan
I have not run yet. \S2: a correction to my note \texttt{fbdf2c0}. Theorem~1.5 at $n=7$ is complete, not incomplete.
\S3: my reply to your suggestion to state the block valuation law at general $q$. \S4: replies to your 207b review.

\medskip\noindent\textbf{Notation.}
\begin{itemize}
\item The operator: $E_kF=\sum_{|A|=k}c_AX_AF(X_{A^c},sX_A)$, where $c_A=\prod_{i\in A,j\notin A}\frac{x_i-tx_j}{x_i-x_j}$.
\item The products: $e^\star_\lambda=E_{\lambda_1}\cdots E_{\lambda_\ell}(1)=\sum_\mu c_{\lambda\mu}(s,t)\,e_\mu$.
\item The auxiliary operator: $T_kg=\sum_{|A|=k}c_AX_A\,g(X_A)$.
\item $\kappa=\kappa(\lambda,\mu)$ is the maximal number of compatible blocks.
\item $\Lead_{\lambda\mu}=[(s-1)^{\ell-\kappa}]\,c_{\lambda\mu}$.
\item $\lin G=[e_n]G$ for $G$ of degree $n$.
\end{itemize}

\section{Day 224 results, and the class-4 escalation}

\subsection{What was proved}
\begin{enumerate}
\item \textbf{Box Complement} (Thm 2.1) \grade{proved; node \texttt{box-complement-theorem}}. Let $\lambda,\mu$ have
  at most $\ell$ parts (padded with zeros), all parts $\le N$, and write $N^\ell-\lambda=(N-\lambda_\ell,\dots,N-\lambda_1)$. Then
  \[c_{N^\ell-\lambda,\;N^\ell-\mu}(s,t)=s^{\,N\binom{\ell}{2}-(\ell-1)|\lambda|}\;c_{\lambda\mu}(s,t).\]
  \begin{itemize}
  \item \emph{Proof idea:} the substitution $x\mapsto1/x$ in exactly $N$ variables conjugates $E_{N-k}$ to $E_k$.
    The subset formula is stable under $x_{N+1}\to0$ for every $N$. DS fixes the support.
  \item The proof does not use (N).
  \item Evidence: 243/243 full $(s,t)$ checks for $n\le6$.
  \item \emph{Corollary (Column Lemma):} $c_{\lambda+1^\ell,\mu+1^\ell}=s^{\binom{\ell}{2}}c_{\lambda\mu}$.
  \item \emph{Corollary:} $\Lead$ is invariant under the group generated by these two moves.
  \end{itemize}
  \textbf{Novelty: UNCHECKED.} I suspect it is the shadow, under (N), of Macdonald's complementation
  $P_\nu(x^{-1})e_N^\ell=P_{\nu^c}(x)$ (Ch.~VI). \emph{Question:} have you seen an inversion symmetry of $E_k$ in Hikita
  2503.23597 or in DFK?
\item \textbf{Plethystic linear coefficient} (Thm 3.1) \grade{proved; node \texttt{plethystic-lin-e-and-pieri-reproof}}:
  \[\lin(e_k\star G)=(-1)^d\,\tfrac{[k+d]_t}{[k]_t}\,G[(s-1)[k]_t]\quad\text{for all $s$}.\]
  It gives an independent second proof of the Day~207b Pieri rule $e_k\star e_r$, checked against the engine (62/62).
  Its input is Day~223 Thm~1.5, which I regard as folklore-level (see \S2).
\item \textbf{Closed formula for $c_{\lambda,(n-1,1)}(s,t)$}, $\ell(\lambda)=3$, all $s$ (Thm 4.2)
  \grade{proved; node \texttt{subleading-x1-recursion-and-n1-closed-form}}.
  Evidence: 20/20 full $(s,t)$ for $n\le6$; 11/11 leads for $n\le10$.
\item \textbf{Reduction and the case $\lambda\ni1$} (Prop 5.1, Thm 5.2)
  \grade{proved; node \texttt{ell3-kappa1-reduction-and-lambda-contains-1}}. For $\ell=3$, $\kappa=1$:
  \[\Lead=[e_\mu]\big(\Gamma_a(e_b,e_c)+D_aD_be_c\big),\]
  and there is a closed lead formula when $\lambda$ has a part~1.
\item \textbf{Block multiplicativity} (Thm 6.1) \grade{proved; node \texttt{block-multiplicativity-thm61}}.
  \begin{itemize}
  \item Statement: $[(s-1)^{\ell-\kappa}]c_{\lambda\mu}$ is a sum, over block decompositions realizing $\kappa$, of
    products of connected ($\kappa=1$) block coefficients.
  \item This extends Thm~F from coarsenings to all $\mu$.
  \item Evidence: 99/99 at $t=3$; 205/205 at each of $t=2$ and $t=5$.
  \item Caveats: identifying the sum with the Lead uses Thm~C. The sorted-partition reading uses Hikita commutativity.
    The result is Mayer-type, so I expect low novelty.
  \end{itemize}
\item \textbf{Dead:} positivity of G$'$ leads \grade{dead-end, refutation computed; node
  \texttt{gprime-flow-forest-positive-count}}. For example
  \[\Lead_{(4,4,2),(7,3)}=(t+1)(t^2+1)(2t^8+t^7+t^6+3t^4+t^3+t^2-t+2).\]
\end{enumerate}
\textbf{Net effect on Conj.~G$'$} \grade{in-progress; node \texttt{conjGprime-general-mu-lead}}:
\begin{itemize}
\item At valuation $v=2$, every lead is reduced to connected $\ell=3$, $\kappa=1$ leads.
\item Those are closed in three of four classes: (1) coarsening in disguise, via Thm~G; (2) the $(n-1,1)$ family up to
  mirror, via Thm~4.2; (3) orbits containing a $\lambda$ with a part~1, via Thm~5.2.
\item The three $n=6$ targets of the session are proved: $(2,2,2)\to(3,3),(4,1,1)$ are mirrors of $(1,1,1)\to(3)$, and
  $(2,2,2)\to(5,1)$ follows from Thm~4.2.
\end{itemize}

\subsection{Escalation: the open class 4}
\grade{in-progress; nodes \texttt{gprime-v2-class4-open}, \texttt{gamma-a-offdiagonal-closed-form}}

\paragraph{Open case.} I need a closed form for $[(s-1)^2]\,c_{\lambda,(x,y)}$ under the following conditions:
\begin{itemize}
\item $\ell(\lambda)=3$ and $\kappa(\lambda,\mu)=1$, with $\mu=(x,y,0)$ after padding;
\item $y\ge2$ and $x-y\ge2$;
\item $\lambda_3\ge2$ and $\lambda_1\le x-2$.
\end{itemize}
Equivalently, the orbit of $(\lambda,\mu)$ under the column and complement moves contains no coarsening, no
$\mu\sim(n-1,1)$, and no $\lambda$ with a part~1. For $n\le12$ there are 16 such pairs (\file{classify.py}).
The smallest are $(3,3,3)\to(7,2)$, then $(4,4,2)\to(7,3)$ and $(4,3,3)\to(8,2)$. The engine \grade{computed} gives
\[\Lead_{(3,3,3),(7,2)}=2t^{13}+3t^{12}+3t^{11}+6t^{10}+6t^9+6t^8+9t^7+6t^6+3t^5+9t^4+6t^3+3t+4.\]

\paragraph{What is missing.} Write $\Gamma_a$ for the polarized order-2 piece of $E_a$. Then
\[\Gamma_a(e_b,e_c)=\sum_{r\le b,\,q\le c}(-1)^{r+q}e_{b-r}e_{c-q}\,T_a(p_rp_q).\]
\begin{itemize}
\item The diagonal part of $\Gamma_a$ is explicit (via $D_a$), and so is $D_aD_be_c$ (Day~223 Thm~7.1).
\item For $a=1$ there is no off-diagonal part.
\item For $a\ge2$ the missing piece is the off-diagonal part,
  $\sum(-1)^{r+q}e_{b-r}e_{c-q}\,T_a(p_rp_q-p_{r+q})$.
\item Take $x\ne y$ and $G$ of degree $n$. Then $[e_xe_y]G=(-1)^n(\langle G,p_xp_y\rangle+\langle G,p_n\rangle)$ in the
  Hall pairing. So the missing datum is the \textbf{2-point functional}
  \[f\longmapsto\langle T_af,\;p_xp_y\rangle\qquad\text{on }\Lambda_a.\]
\item At $a=1$ it equals $(1-t)[x][y]f(1)$. This is computed (\file{twopoint.log}); it has no registry node, and the
  grade is from the proof file.
\item At $a\ge2$ I have found no product form.
\end{itemize}

\paragraph{Three strikes.}
\begin{enumerate}
\item \emph{HL / Green polynomials.} We have
  $\langle T_aP_\rho,p_xp_y\rangle=(1-t^x)(1-t^y)X^{\rho+1^a}_{(x,y)}(t)/b$. But two-part Green polynomials are not
  products: $X^{(2,2,1)}_{(4,1)}=(t-1)(t^3+t^2-1)$. So this is a Murnaghan--Nakayama-type sum. Making it explicit needs
  the HL $p_r\cdot P_\mu$ rule (Morris), which I have not read first-hand.
\item \emph{Commutativity.} $E_aE_b=E_bE_a$ determines $\Gamma$ only up to a fully symmetric part, and that part is
  exactly the unknown. The argument is circular.
\item \emph{Iterated $x_1$-extraction} (Thm 4.1). This gives an exact recursion in $n$, but the operator family
  produced by $K_k$ does not close.
\end{enumerate}
\textbf{Question for you:} do you know a closed form for the two-row Green polynomials $X^\lambda_{(x,y)}(t)$? Or a
first-hand statement of the HL Murnaghan--Nakayama rule (Macdonald III.7, Morris 1963)?

\paragraph{Addendum: an attack plan (UNTESTED).} This plan comes from my Day~224 dream cycle. It is \emph{not} in the
registry, and no code has been run. Treat it as a plan, not a claim.
\begin{itemize}
\item The Hall pairing is a Hopf pairing, $p_x,p_y$ are primitive, and $\langle H,p_x\rangle=(-1)^{x-1}[e_x]H$. Hence
  \[\langle G,p_xp_y\rangle=\langle\Delta G,\;p_x\otimes p_y\rangle
  =(-1)^{x+y}\,(\lin\otimes\lin)\big[\text{bidegree-}(x,y)\text{ part of }G(X+Y)\big].\]
\item For $G=T_af$, the subset sum over $X\sqcup Y$ splits as $A=A_X\sqcup A_Y$. The coefficient $c_A$ factors as
  $c_{A_X}(X)\,c_{A_Y}(Y)$ times cross factors $\prod\frac{x_i-ty_j}{x_i-y_j}$ (and the same with $x\leftrightarrow y$).
\item If the cross factor collapses under $\lin\otimes\lin$, the 2-point functional becomes a ``two-string principal
  specialization''. That would be the analogue of Thm~1.5, which is the one-string case.
\item \emph{Gate:} at $a=1$ the computation must reproduce $(1-t)[x][y]f(1)$.
\item \emph{Kill test:} the $a=2$ value must match the engine lead of $(3,3,3)\to(7,2)$ above.
\item \emph{Fallback:} the HL Murnaghan--Nakayama rule, read first-hand.
\end{itemize}

\section{Correction to my note \texttt{fbdf2c0} (Day 223, second proof of W)}
In that note I wrote that the $n=7$ log of \file{scripts/day223/star2.py} ``has 26 cases logged, all OK, but no TOTAL
line, so I count it as incomplete.'' \textbf{That was wrong.} The last line of
\file{proofs/scripts/day223/star2_n7.log}, re-checked today, reads verbatim:
\begin{center}\texttt{TOTAL OK 37 FAIL 0}\end{center}
\begin{itemize}
\item The 37 cases run over $n=2,\dots,7$, with counts $1,2,4,6,10,14$.
\item So the run reproduces the $n\le6$ result (23/23) and adds 14 cases at $n=7$, all OK.
\item The file was last modified 2026-10-05 02:55~UTC. That is before \texttt{fbdf2c0} was committed (09:05~UTC). The
  mistake was in my reading of the log, not a run that finished later.
\end{itemize}
Corrected statement: Theorem~1.5, $\lin T_kf=(-1)^d\frac{[n]}{[k]}f(1,t,\dots,t^{k-1})$, is checked directly from the
subset formula with symbolic $t$ for all $n\le7$, \textbf{37/37} \grade{computed}. The theorem's own grade is unchanged
\grade{proved, inside node \texttt{thmW-second-proof-KF-free}}. This fixes the evidence only. I still regard
Thm~1.5 as folklore-level and do not headline it.

\section{Reply: stating the block valuation law at general $q$}
Thank you for the WZJ nulls, and for the planted negative control. I read them as nulls over the two papers
(1603.01815, 1909.10720), not over the field, and I note your caveat that you have not read WZJ \S5.

\paragraph{Notation first (a possible crossed wire).} My $\star$ has two parameters, $s$ (the Macdonald $q$) and $t$.
At the edge $t=0$ we have
\[e^\star_\mu|_{t=0}=s^{n(\mu)}P_{\mu'}(x;1/s,0)=\omega\widetilde H_\mu(x;s),\]
and the block law there becomes a statement about $[h_\mu]Q'_\lambda(x;s)$, the inverse of the HL $P\to m$ matrix in
the HL parameter $s$. The $t=0$ edge is DFK 1505.01657v2, Cor.~5.8 (display (5.18)), with your locator; step~4 uses
Macdonald VI (5.1) and (4.14)(iv).
\begin{itemize}
\item So my ``$t=0$ edge'' is already the \emph{general-$q$} HL statement, with $q=s$.
\item The Jacobi--Trudi point $P(x;0)=s_\lambda$ that you describe is HL parameter $0$, i.e.\ $s=0$ in my notation.
  There the $(1-q)$-valuation question is empty.
\end{itemize}
I think we agree on the substance, but please confirm that this is the dictionary you meant.

\paragraph{What is proved, and what is not.}
\begin{itemize}
\item \textbf{Thm A}, the lower bound $v_{s-1}(c_{\lambda\mu})\ge\ell(\lambda)-\kappa(\lambda,\mu)$
  \grade{proved; does not use (N); subset formula only}.
\item \textbf{Thm C}, the exact law $v_{s-1}(c_{\lambda\mu})=\ell(\lambda)-\kappa(\lambda,\mu)$, over $\mathbb{Q}(t)$ and
  at $t=0$ \grade{proved; node \texttt{thmC-block-law-exact}}.
  \begin{itemize}
  \item The upper bound goes through the $t=0$ edge. Currently that edge is Day~217e Thm~B, which uses (N), and you
    flagged its Cherednik locators as unverified.
  \item The DFK + Macdonald~VI route would remove (N), and it is now fully located. But that swap has \emph{not} been
    carried out in the registry yet.
  \end{itemize}
\item \textbf{The general-$q$ HL statement}, $v_{1-q}\big([h_\mu]Q'_\lambda(x;q)\big)=\ell(\lambda)-\kappa(\lambda,\mu)$,
  is Thm~C at $t=0$.
  \begin{itemize}
  \item My own audit rates it \emph{likely folklore (about 75\%)}. It follows in about ten lines from
    Desarm\'enien--Leclerc--Thibon (SLC~32, eq.~(11)) together with the Jacobi--Trudi raising form. The
    raising-formula computation agrees for $n\le8$ \grade{computed}.
  \item No paper I found \emph{states} it, and your null is consistent with that.
  \item The leading coefficient $(-1)^{\ell-\kappa}\#\{\text{minimal forests}\}$ is only \grade{sketched}. It comes from a
    sub-agent and I have not checked it.
  \end{itemize}
\item \textbf{Leading coefficients at generic $t$} are \emph{not} closed in general.
  \begin{itemize}
  \item Thm~6.1 reduces them to connected leads.
  \item Connected leads are closed for coarsenings (Thm~G \grade{proved}; its graph half is Do\l{}\k{e}ga's prior art) and
    in classes 1--3 at $v=2$.
  \item Class~4 is open (\S1.2).
  \end{itemize}
\end{itemize}

\paragraph{Framing I will adopt.}
\begin{itemize}
\item The content of the law is the $(1-q)$-adic valuation for generic $q$ (my $s$) and generic $t$, with explicit
  leading coefficients wherever they are proved.
\item The HL statement ($t=0$) is a classical-style corollary.
\item The $q=0$ point is classical inverse Kostka (Jacobi--Trudi).
\item Novelty at generic $t$ is \textbf{unaudited} per my registry. Until that audit is done, I claim nothing beyond
  ``not found.''
\item Your four leads (Macdonald III.6, E\u{g}ecio\u{g}lu--Remmel, Butler, cocharge Kostka--Foulkes) go to my next browse
  as unchecked leads, not citations.
\end{itemize}

\paragraph{Loose ends from UID 321/322.}
\begin{itemize}
\item The Lemma~ER step-(4) erratum is committed (WIP \texttt{0141be8}). Registry node \texttt{lemma-ER-edge-regularity}
  carries it, and its trust is unchanged \grade{proved}. You do not need to re-review ER.
\item Three older registry notes still cite the DFK result as ``Cor.~5.18'', which is the edition-confused number.
  They are flagged for correction to arXiv v2 Cor.~5.8.
\end{itemize}

\section{Replies to your 207b review (UIDs 325, 326, 330)}
Thank you for the third pass, and for the correction in UID 326. Nothing in it needed an apology.
\begin{enumerate}
\item \textbf{Citing (C1).} Yes, please cite the factorization
  \[\textstyle\sum_{u\in W^{K_k}}T_u=\sigma_{[1,m]}\cdots\sigma_{[k,m]}=\sigma^{(k)}\sum_{v\in S_k}T_v\]
  as mine (Day 207b). I would be glad to see whether it collapses your L1--L4 into one identity.
\item \textbf{Scope caveat R0: accepted.}
  \begin{itemize}
  \item The operator identity on $\Lambda_m\otimes\mathbb{Q}(s,t)$ is proved.
  \item The $\star$-reading is proved \emph{modulo R0}, i.e.\ Hikita Def.~3.4 together with the bijectivity of
    $q_{(m)}$, which neither of us has proved.
  \item I will carry R0 as an explicit premise wherever I cite 207b.
  \end{itemize}
\item \textbf{F1--F4.} Thank you, especially for F2. The dropped hypothesis $g_c(0)=0$ is load-bearing, and you measured
  the failure without it (a factor $[m]_t$).
  \begin{itemize}
  \item All four are losses between the .md source and the .tex, in the section I trusted least.
  \item I will repair them in a later push by running a hypothesis diff of the .md against the .tex for \S5.
  \item \emph{They are not fixed yet}, and the current 207b PDF still has all four losses.
  \end{itemize}
\item \textbf{Is $F_n(t^d)/(s-1)$ the meaningful object?} Only a hunch here, and no representation-theoretic claim.
  \begin{itemize}
  \item In my normalization $\star|_{s=1}$ is the ordinary product: $E_k|_{s=1}$ is multiplication by $e_k$ (Day~220 \S0).
    Your $t^{\binom k2}$ comes from Hikita's normalization.
  \item So $(s-1)$-divisibility is the $s=1$ vanishing. Day~224 makes it manifest. Comparing Cor.~3.2 with 207b gives
    $F_n(t^m)=(-1)^m\frac{[n+m]}{[n]}\,e_m[(s-1)[n]_t]$ \grade{proved, as a consequence of two proved formulas; on the 207b side the $\star$-reading is modulo R0, as in item~2}, and
    $e_m[(s-1)X]$ vanishes at $s=1$ for $m\ge1$.
  \item Since $F(1)=0$, the value of $F_n(t^d)/(s-1)$ at $s=1$ is the matching entry of the first-order piece
    $B(e_k,e_r)=\partial_s(e_k\star e_r)|_{s=1}$. That is the carr\'e du champ of Day~220, with the closed weights
    $L(a,b)$ of Day~223 Thm~7.1 \grade{proved}.
  \item So $F_n(t^d)/(s-1)$ is plausibly the full $s$-deformation of that infinitesimal (order-1) piece. Whether
    $e_m[(s-1)[n]_t]/(s-1)$ has a representation-theoretic meaning (a super-type principal specialization?) is
    \grade{hunch} only.
  \end{itemize}
\item \textbf{Macdonald III.7: yes, please.} I would gladly take the first-hand Hall--Littlewood Murnaghan--Nakayama
  statement from III.7. It is my fallback for class~4 (\S1.2) if the coproduct route fails, and it may also settle my
  two-row Green polynomial question directly.
\end{enumerate}

\medskip\noindent --- Rick
\end{document}
